% ===============================================================================
% SL-04d-ML: Stundenleistung MUSTERLÖSUNG
% Spanisch Klasse 7 - Sequenz 4: Mi colegio
% ===============================================================================

\documentclass[11pt, a4paper]{article}

\newif\ifswmodus
\swmodusfalse

\usepackage[utf8]{inputenc}
\usepackage[T1]{fontenc}
\usepackage[ngerman]{babel}
\usepackage{helvet}
\renewcommand{\familydefault}{\sfdefault}

\usepackage[margin=2cm]{geometry}
\usepackage{enumitem}
\usepackage{tabularx}
\usepackage{booktabs}

\usepackage{xcolor}
\usepackage{tcolorbox}
\tcbuselibrary{skins}

\usepackage{fancyhdr}
\usepackage{lastpage}
\usepackage{needspace}

\definecolor{spanisch-color}{HTML}{c41e3a}
\definecolor{grau-color}{HTML}{888888}
\definecolor{hellgrau-color}{HTML}{f5f5f5}
\definecolor{loesung-color}{HTML}{c62828}

\definecolor{sw-dark}{HTML}{000000}
\definecolor{sw-gray}{HTML}{666666}
\definecolor{sw-white}{HTML}{ffffff}

\ifswmodus
    \colorlet{spanisch}{sw-dark}
    \colorlet{grau}{sw-gray}
    \colorlet{hellgrau}{sw-white}
    \colorlet{loesung}{sw-dark}
\else
    \colorlet{spanisch}{spanisch-color}
    \colorlet{grau}{grau-color}
    \colorlet{hellgrau}{hellgrau-color}
    \colorlet{loesung}{loesung-color}
\fi

\newcommand{\fachfarbe}{spanisch}

\newcommand{\thema}{Stundenleistung: Mi colegio}
\newcommand{\fach}{Spanisch}
\newcommand{\klasse}{7}
\newcommand{\maxpunkte}{20}
\newcommand{\zeit}{25 Minuten}
\newcommand{\datum}{\today}

\pagestyle{fancy}
\fancyhf{}
\renewcommand{\headrulewidth}{0.4pt}
\renewcommand{\footrulewidth}{0.4pt}

\lhead{\fach{} -- Klasse \klasse}
\chead{\textcolor{loesung}{\textbf{SL-04d MUSTERLÖSUNG}}}
\rhead{\datum}

\lfoot{\zeit}
\cfoot{}
\rfoot{Seite \thepage\ von \pageref{LastPage}}

\newcommand{\punktebox}[1]{\hfill\fbox{\small \_/#1}}
\newcommand{\luecke}[1]{\rule{#1}{0.4pt}}
\newcommand{\lsg}[1]{\textcolor{loesung}{\textbf{#1}}}

\newtcolorbox{infobox}{
    colback=hellgrau,
    colframe=\fachfarbe,
    boxrule=1pt,
    arc=0pt,
    left=8pt, right=8pt, top=8pt, bottom=8pt
}

\newenvironment{aufgabe}[1]{%
    \needspace{5\baselineskip}%
    \nopagebreak[4]%
    \vspace{10pt}
    \noindent\textbf{\textcolor{\fachfarbe}{Aufgabe #1}}
    \nopagebreak[4]%
    \vspace{5pt}
    \nopagebreak[4]%
    \begin{list}{}{\leftmargin=0pt}
    \item[]
}{%
    \end{list}
}

\newcommand{\es}[1]{\textcolor{\fachfarbe}{\textbf{#1}}}

\begin{document}

\begin{center}
    {\Large\bfseries\textcolor{\fachfarbe}{\thema}}\\[0.3em]
    {\large\textcolor{loesung}{\textbf{--- MUSTERLÖSUNG ---}}}
\end{center}

\vspace{5pt}

\noindent
\begin{tabularx}{\textwidth}{|X|X|X|}
\hline
\textbf{Name:} \lsg{Musterschüler/in} &
\textbf{Klasse:} \klasse &
\textbf{Datum:} \datum \\
\hline
\end{tabularx}

\vspace{10pt}

\begin{infobox}
\textbf{Zeit:} \zeit{} | \textbf{Punkte:} \maxpunkte{} | \textbf{Hilfsmittel:} keine
\end{infobox}

\vspace{15pt}

% ===============================================================================
% AUFGABE 1: Schulfächer zuordnen (4P)
% ===============================================================================

\begin{aufgabe}{1 -- Schulfächer zuordnen (Asignaturas)}
Verbinde das spanische Schulfach mit der deutschen Bedeutung. \punktebox{4}
\end{aufgabe}

\begin{center}
\begin{tabular}{rcl}
\es{el inglés} & \lsg{---} & Mathematik \lsg{$\rightarrow$ Englisch} \\[0.8em]
\es{las matemáticas} & \lsg{---} & Musik \lsg{$\rightarrow$ Mathematik} \\[0.8em]
\es{la música} & \lsg{---} & Sport \lsg{$\rightarrow$ Musik} \\[0.8em]
\es{la educación física} & \lsg{---} & Englisch \lsg{$\rightarrow$ Sport} \\
\end{tabular}
\end{center}

\textcolor{loesung}{\textit{Je 1P pro richtige Zuordnung = 4P}}

\vspace{15pt}

% ===============================================================================
% AUFGABE 2: Wochentage + Uhrzeiten (4P)
% ===============================================================================

\begin{aufgabe}{2 -- Wochentage und Uhrzeiten}
Ergänze die fehlenden Wörter. \punktebox{4}
\end{aufgabe}

\noindent
\textbf{a) Ergänze die Wochentage:} (2P)

\vspace{0.5em}

\noindent
lunes, martes, \lsg{miércoles}, jueves, \lsg{viernes}

\textcolor{loesung}{\textit{Je 1P pro richtigem Wochentag = 2P}}

\vspace{1em}

\noindent
\textbf{b) Übersetze die Uhrzeiten ins Spanische:} (2P)

\vspace{0.5em}

\noindent
Um 9 Uhr = \lsg{A las nueve}

\vspace{0.8em}

\noindent
Um halb 11 = \lsg{A las diez y media}

\textcolor{loesung}{\textit{Je 1P pro richtige Uhrzeit = 2P}}

\vspace{15pt}

% ===============================================================================
% AUFGABE 3: -ar Verben konjugieren (6P)
% ===============================================================================

\begin{aufgabe}{3 -- Verben auf -ar konjugieren}
Setze die richtige Form des Verbs ein. \punktebox{6}
\end{aufgabe}

\noindent
a) Yo \lsg{estudio} (estudiar) español todos los días. \hfill (1P)

\vspace{0.8em}

\noindent
b) Tú \lsg{hablas} (hablar) muy bien inglés. \hfill (1P)

\vspace{0.8em}

\noindent
c) María \lsg{canta} (cantar) en el coro de la escuela. \hfill (1P)

\vspace{0.8em}

\noindent
d) Nosotros \lsg{escuchamos} (escuchar) música en clase de música. \hfill (1P)

\vspace{0.8em}

\noindent
e) Las clases \lsg{terminan} (terminar) a las dos. \hfill (1P)

\vspace{0.8em}

\noindent
f) Yo \lsg{bailo} (bailar) salsa los viernes. \hfill (1P)

\textcolor{loesung}{\textit{Je 1P pro richtige Verbform = 6P}}

\vspace{15pt}

% ===============================================================================
% AUFGABE 4: Schulalltag beschreiben (6P)
% ===============================================================================

\begin{aufgabe}{4 -- Mi día en el colegio}
Schreibe 3 Sätze über deinen Schulalltag. Verwende Schulfächer, Wochentage und/oder Uhrzeiten. \punktebox{6}
\end{aufgabe}

\noindent
\textbf{Beispiel:} Los lunes tengo matemáticas a las ocho.

\vspace{0.8em}

\noindent
1. \lsg{Los martes tengo español a las nueve.} \hfill (2P)

\vspace{1em}

\noindent
2. \lsg{Estudio historia los miércoles a las diez.} \hfill (2P)

\vspace{1em}

\noindent
3. \lsg{La educación física es los viernes.} \hfill (2P)

\vspace{0.5em}

\textcolor{loesung}{\textit{Bewertung pro Satz (max. 2P):}}
\begin{itemize}[leftmargin=2em]
    \item \textcolor{loesung}{\textit{2P = korrekter Satz mit Schulfach + Wochentag oder Uhrzeit}}
    \item \textcolor{loesung}{\textit{1P = Satz mit kleinen Fehlern oder nur einem Element}}
    \item \textcolor{loesung}{\textit{0P = unverständlich oder komplett falsch}}
\end{itemize}

% ===============================================================================
% PUNKTEÜBERSICHT
% ===============================================================================

\vspace{25pt}
\hrule
\vspace{10pt}

\noindent
\begin{tabularx}{\textwidth}{|X|c|c|c|c||c|}
\hline
\textbf{Aufgabe} & 1 & 2 & 3 & 4 & \textbf{Gesamt} \\
\hline
\textbf{max. Punkte} & 4 & 4 & 6 & 6 & \textbf{20} \\
\hline
\textbf{erreicht} & \lsg{4} & \lsg{4} & \lsg{6} & \lsg{6} & \lsg{20} \\
\hline
\end{tabularx}

\vspace{10pt}

\begin{flushright}
\textbf{Note:} \lsg{1}
\end{flushright}

\vspace{1em}

\textbf{Notenschlüssel:}

\begin{tabular}{|c|c|c|c|c|c|c|}
\hline
\textbf{Note} & 1 & 2 & 3 & 4 & 5 & 6 \\
\hline
\textbf{ab Punkte} & 19 & 16 & 12 & 8 & 4 & <4 \\
\hline
\end{tabular}

\end{document}
